\chapter{C++ pour la latence III~: le compilateur}\label{ch:ll:cpp-for-latency-iii-the-compiler}

En juillet 2009, un pilote du noyau Linux a été trouvé contenant les lignes \texttt{struct sock *sk = tun->sk;} et, un peu
plus bas, \texttt{if (!tun) return POLLERR;}. Le test devait protéger contre un pointeur nul. Le compilateur avait
supprimé~: le pointeur avait déjà été déréférencé, or un déréférencement nul est un \omterm{def:ll:cpp-for-latency-iii-the-compiler:ub}{comportement indéfini}, donc le pointeur ne pouvait pas être nul et le test était
du code mort. Sur une machine où l'adresse zéro pouvait être mappée, la suppression est devenue un exploit. Le
compilateur n'a rien fait de mal~; il a lu le contrat du programme à la lettre. Ce chapitre porte sur ce contrat et sur les
autres leviers que le compilateur offre à l'ingénieur de latence~: les niveaux d'optimisation et les options de cible, l'\omterm{def:ll:cpp-for-latency-iii-the-compiler:pgo}{optimisation guidée par profil}
et l'\omterm{def:ll:cpp-for-latency-iii-the-compiler:lto}{optimisation à l'édition de liens}, et l'habitude qui relie le tout, lire l'assembleur et mesurer chaque option contre une somme de contrôle.

\section{Comportement indéfini comme contrat d'optimisation}

\begin{definition}[Comportement indéfini]\label{def:ll:cpp-for-latency-iii-the-compiler:ub}
Un \emph{comportement indéfini}\index{comportement indéfini} est un comportement pour lequel la norme C++ n'impose aucune exigence~:
un dépassement de capacité d'un entier signé, un déréférencement nul ou pendant, un accès hors des bornes, une course aux données, la lecture d'un objet à travers
un pointeur d'un type sans rapport. Le compilateur peut supposer que cela ne se produit jamais, et optimise chaque programme comme si cela ne se produisait pas.
\end{definition}

L'hypothèse est à l'origine d'une bonne part de la puissance de l'optimiseur. Comme un dépassement d'entier signé ne peut pas se produire, \texttt{x + 1 > x}
est vrai, un compteur de boucle ne peut pas revenir à zéro, et le nombre d'itérations d'une boucle peut être calculé à l'avance~; comme le déréférencement d'un pointeur nul ne peut pas
se produire, un pointeur déréférencé une fois est non nul ensuite. Dans les
\cref{lst:ll:cpp-for-latency-iii-the-compiler:plus,lst:ll:cpp-for-latency-iii-the-compiler:null} montrent les deux cas sur le compilateur de ce livre
la comparaison se compile en «~return 1~» même sans optimisation, et le test de nullité du motif de 2009 disparaît
à \texttt{-O2}.

\omcode{code/low-latency/08-cpp-for-latency-iii-the-compiler/cpp/asm/plus_one_greater_O2.s}{1}{2}{\texttt{plus\_one\_greater(int x)}, \texttt{return x + 1 > x;}, à -O2 (g++ 11.4)~: une constante.}\label{lst:ll:cpp-for-latency-iii-the-compiler:plus}
\omcode{code/low-latency/08-cpp-for-latency-iii-the-compiler/cpp/asm/first_then_check_O2.s}{1}{2}{\texttt{first\_then\_check(const int* p)}~: déréférencement, puis test de nullité~; le test a disparu.}\label{lst:ll:cpp-for-latency-iii-the-compiler:null}

Pour un système de trading, la leçon est double. Le \omterm{def:ll:cpp-for-latency-iii-the-compiler:ub}{comportement indéfini} n'est pas «~la machine fait quelque chose de bizarre~»~; c'est une autorisation
donnée au compilateur de supprimer des tests, de réordonner des travaux et de supposer des invariants que le code viole, et le résultat peut changer
avec la version du compilateur ou avec une option. Or ce même contrat est ce qui rend le C++ rapide~: les idiomes qui donnent à l'optimiseur
des faits (compteurs de boucle signés, absence d'aliasing entre types sans rapport) sont ceux qui lui permettent de garder les valeurs dans les registres et
de vectoriser. Les outils qui détectent les violations s'exécutent dans les constructions de test~: le sanitiseur de \omterm{def:ll:cpp-for-latency-iii-the-compiler:ub}{comportement indéfini}
(\texttt{-fsanitize=undefined}) signale chaque dépassement au moment où il se produit, le sanitiseur d'adresses chaque accès invalide.

\section{Aliasing}

\begin{definition}[Aliasing strict]\label{def:ll:cpp-for-latency-iii-the-compiler:alias}
Selon la règle de l'\emph{aliasing strict}\index{aliasing strict}, un objet ne peut être accédé qu'à travers une lvalue de son propre
type (ou d'un type très proche~: une variante signée ou non signée, un type caractère), de sorte que le compilateur peut supposer que des pointeurs
vers des types sans rapport ne désignent jamais la même mémoire. GCC applique la règle à partir de \texttt{-O2}.
\end{definition}

\omcode{code/low-latency/08-cpp-for-latency-iii-the-compiler/cpp/asm/alias_O2.s}{1}{4}{\texttt{alias(int* i, float* f)}~: écrit 1 par \texttt{i}, 0 par \texttt{f}, renvoie \texttt{*i}~; le compilateur renvoie 1 sans relire.}\label{lst:ll:cpp-for-latency-iii-the-compiler:alias}

Dans le \cref{lst:ll:cpp-for-latency-iii-the-compiler:alias}, le compilateur renvoie 1 sans relire \texttt{*i}, parce qu'une
écriture par un \texttt{float*} ne peut pas modifier un \texttt{int}. Appelée avec deux pointeurs vers les mêmes quatre octets, la fonction
renvoie 1 alors que la mémoire contient 0. Les décodeurs sont les premiers touchés~: lire un prix dans un tampon d'octets en convertissant
l'adresse du tampon en \texttt{std::int32\_t*} constitue une violation d'aliasing et, pour les champs non alignés, aussi un \omterm{def:ll:cpp-for-latency-iii-the-compiler:ub}{comportement indéfini}. Les formes correctes
\texttt{std::memcpy} vers une variable locale ou \texttt{std::bit\_cast} de C++20 se compilent en un seul chargement~; le compilateur sait
ce qu'elles font (le chapitre~16 construit ses codecs de cette façon).

\section{Options et choix de la cible}

Le niveau d'optimisation est un ensemble de passes~: \texttt{-O2} est le niveau de publication conventionnel, \texttt{-O3} ajoute
un \omterm{def:ll:cpp-for-latency-ii-compile-time:inline}{inlining} plus agressif, un clonage et une vectorisation, et \texttt{-Ofast} ajoute \texttt{-ffast-math}. \texttt{-march=native}
permet au compilateur d'utiliser chaque instruction de la machine de compilation (AVX2 sur ce portable), ce qui rend le binaire inutilisable sur des machines plus
anciennes~; les constructions de production nomment la cible explicitement. \texttt{-ffast-math} autorise le compilateur à réordonner
les opérations en virgule flottante comme si elles étaient associatives et à supposer qu'il n'existe ni NaN ni infini.

\begin{definition}[Multiversionnage de fonctions]\label{def:ll:cpp-for-latency-iii-the-compiler:fmv}
Le \emph{multiversionnage de fonctions}\index{multiversionnage de fonctions} compile une même fonction plusieurs fois pour différents
jeux d'instructions et choisit une version à l'exécution selon les fonctionnalités du processeur~: l'attribut de GCC
\texttt{target\_clones(\textquotedbl{}avx2\textquotedbl{}, \textquotedbl{}default\textquotedbl{})} construit les deux versions et répartit l'appel au chargement du programme.
\end{definition}

La \Cref{fig:ll:cpp-for-latency-iii-the-compiler:flags} mesure le décodeur de flux du livre~1 avec un gestionnaire par message
dans une seconde unité de traduction, sous sept jeux d'options, et une somme en virgule flottante sous trois. Sur le décodeur,
\texttt{-O3} a été plus de trois fois plus rapide que \texttt{-O2}, l'\omterm{def:ll:cpp-for-latency-iii-the-compiler:pgo}{optimisation guidée par profil} seule deux fois et demie,
et la combinaison de tout près de cinq fois~; \texttt{-march=native} et l'\omterm{def:ll:cpp-for-latency-iii-the-compiler:lto}{optimisation à l'édition de liens} seules ne changent presque
rien. Sur la somme flottante, \texttt{-ffast-math} a été sept fois plus rapide et a donné un autre résultat, car elle a sommé dans
un autre ordre~: la somme de contrôle du benchmark l'a détecté, et \texttt{firm.flagbench} ne classe pas un jeu d'options dont la sortie
diffère.

\begin{omfigure}
\begin{tikzpicture}
\begin{axis}[name=dec,width=8.2cm,height=6cm,xbar,bar width=7pt,xmin=0,xmax=5,
  symbolic y coords={O0,O2,O2 native,O2 LTO,O2 PGO,O3,O3 native LTO PGO},ytick=data,enlarge y limits=0.1,
  xlabel={accélération par rapport à -O2},title={\small décodeur de flux},tick label style={font=\footnotesize},label style={font=\small},
  grid=major,grid style={gray!25},nodes near coords,nodes near coords style={font=\scriptsize},
  every node near coord/.append style={/pgf/number format/fixed,/pgf/number format/precision=2},table/col sep=comma]
\addplot[fill=omDef!55,draw=omDef] table[y=name,x=speedup]{figdata/low-latency/08-cpp-for-latency-iii-the-compiler/measured_flags.csv};
\end{axis}
\begin{axis}[at={(dec.east)},anchor=west,xshift=1.9cm,width=5.8cm,height=6cm,xbar,bar width=9pt,xmin=0,xmax=9,
  symbolic y coords={O2,O3 native,O3 native fast-math},ytick=data,yticklabels={O2,O3 native,{O3 native\\fast-math}},
  y tick label style={align=right},enlarge y limits=0.3,xlabel={accélération par rapport à -O2},title={\small somme de flottants},
  tick label style={font=\footnotesize},label style={font=\small},grid=major,grid style={gray!25},nodes near coords,
  nodes near coords style={font=\scriptsize},every node near coord/.append style={/pgf/number format/fixed,/pgf/number format/precision=1},
  table/col sep=comma]
\addplot[fill=omThm!45,draw=omThm] table[y=name,x=speedup]{figdata/low-latency/08-cpp-for-latency-iii-the-compiler/measured_fastmath.csv};
\end{axis}
\end{tikzpicture}

\omcaption{À gauche~: le benchmark de décodage (échantillon de 2\,020 messages du livre~1 décodé 300 fois, avec un gestionnaire par message dans une autre
unité de traduction) sous sept jeux d'options~; chaque jeu a donné la même somme de contrôle. À droite~: somme de $2^{20}$ flottants~; avec
\texttt{-ffast-math}, la somme a changé. Meilleur de cinq exécutions, g++ 11.4. Mesures sur un portable (Intel Core Ultra 7 155H) sous
WSL2, sans cœurs isolés, épinglé à un seul processeur. Données~: \texttt{bench\_flags.py}.}\label{fig:ll:cpp-for-latency-iii-the-compiler:flags}
\end{omfigure}

\section{Optimisation guidée par profil et optimisation à l'édition de liens}

\begin{definition}[Optimisation guidée par profil]\label{def:ll:cpp-for-latency-iii-the-compiler:pgo}
Dans l'\emph{optimisation guidée par profil}\index{optimisation guidée par profil} (PGO), le programme est compilé avec une instrumentation,
exécuté sur une charge de travail représentative qui enregistre la fréquence de chaque sens de chaque branche et les valeurs et appels les plus courants,
et recompilé avec ce profil, de sorte que le compilateur insère, déroule et dispose le code selon le comportement mesuré
plutôt que selon des suppositions.
\end{definition}

\begin{definition}[Optimisation à l'édition de liens]\label{def:ll:cpp-for-latency-iii-the-compiler:lto}
Dans l'\emph{optimisation à l'édition de liens}\index{optimisation à l'édition de liens} (LTO), le compilateur conserve sa représentation intermédiaire
dans les fichiers objets et réoptimise l'ensemble du programme au moment de l'édition de liens, comme si toutes les unités de traduction n'en formaient qu'une~: des fonctions
peuvent être insérées d'un fichier à l'autre.
\end{definition}

\begin{omfigure}
\begin{tikzpicture}[font=\footnotesize,
  st/.style={draw=omDef,fill=omDef!8,rounded corners=2pt,minimum height=0.75cm,align=center,text width=2.1cm},
  io/.style={draw=gray,fill=gray!10,rounded corners=2pt,minimum height=0.75cm,align=center,text width=1.9cm}]
\node[st] (a) at (0,0) {construction\\ instrumentée};
\node[st] (b) at (3.1,0) {entraîner sur\\ un jour rejoué};
\node[io] (c) at (6.0,0) {profil\\ (comptages)};
\node[st] (d) at (8.9,0) {reconstruire\\ avec le profil};
\node[io] (e) at (11.6,0) {binaire de\\ production};
\foreach \x/\y in {a/b,b/c,c/d,d/e} \draw[-Stealth,thick] (\x) -- (\y);
\node[st] (f) at (0,-1.6) {chaque fichier\\ en code\\ intermédiaire};
\node[st,text width=2.7cm] (g) at (3.1,-1.6) {édition de liens~:\\ tout optimiser\\ comme un seul};
\node[io] (h) at (6.0,-1.6) {binaire};
\draw[-Stealth,thick] (f) -- (g); \draw[-Stealth,thick] (g) -- (h);
\node[anchor=east,font=\footnotesize\bfseries] at (-1.3,0) {PGO};
\node[anchor=east,font=\footnotesize\bfseries] at (-1.3,-1.6) {LTO};
\end{tikzpicture}

\omcaption{L'\omterm{def:ll:cpp-for-latency-iii-the-compiler:pgo}{optimisation guidée par profil} (PGO) construit deux fois autour d'une exécution d'entraînement, dont la charge de travail détermine ce pour quoi le compilateur
optimise~; l'\omterm{def:ll:cpp-for-latency-iii-the-compiler:lto}{optimisation à l'édition de liens} (LTO) reporte l'optimisation de l'ensemble du programme à l'édition de liens.}\label{fig:ll:cpp-for-latency-iii-the-compiler:pgo}
\end{omfigure}

L'exécution d'entraînement est l'entrée critique de la PGO. Un profil relevé un matin calme optimise les chemins du matin calme~; un
système de trading doit être entraîné sur un rejeu d'une journée chargée et représentative (chapitre~23), et réentraîné lorsque la charge de travail
change, faute de quoi il sera optimisé pour les mauvaises branches. Le bénéfice de la LTO dépend de la structure du code~: ici, le
gestionnaire de chaque message était petit et son coût était dominé par les appels propres au décodeur, si bien que la LTO seule n'a rien gagné.
Les deux allongent la compilation et rendent le binaire tributaire de davantage que la source~; les deux ont leur place dans le pipeline
de publication, mesurées par la matrice d'options à chaque version, et non appliquées une fois pour toutes puis oubliées.

\section{Lire l'assembleur}

La sortie du compilateur répond à des questions qu'aucun benchmark ne peut trancher~: si une fonction a été mise en ligne, si une boucle a été vectorisée,
si un contrôle a disparu, si une valeur est conservée dans un registre. \texttt{g++ -O2 -S -masm=intel} l'écrit pour un fichier~;
\texttt{objdump -d -M intel} désassemble un binaire~; les explorateurs de compilateurs en ligne la montrent ligne à ligne. Les listings d'assembleur du livre sont
générés par des scripts à partir de sources testées, et leurs tests vérifient des propriétés (aucun appel, un retour constant) avec le compilateur installé,
jamais les octets exacts, qui changent avec le compilateur.

\begin{method}[Choisir les options de publication]\label{met:ll:cpp-for-latency-iii-the-compiler:flags}
\begin{enumerate}
\item Donner au benchmark une somme de contrôle de sa sortie et une charge de travail réaliste (un jour rejoué, et non un message répété dans une boucle).
\item Compiler le benchmark sous une matrice d'ensembles d'options, dont des ensembles guidés par profil entraînés sur un autre jour que celui mesuré.
\item Rejeter tout ensemble dont la somme de contrôle diffère de celle de la version de référence, quelle que soit sa vitesse.
\item Classer les autres d'après les centiles mesurés sur des exécutions répétées (chapitre~25), et lire l'assembleur des fonctions
critiques pour comprendre le gagnant.
\item Figer les options, la version du compilateur et la cible dans le système de construction~; relancer la matrice à chaque mise à niveau du compilateur.
\end{enumerate}
\end{method}

\section{Tutoriel~: une matrice d'options et trois contrats}\label{tut:ll:cpp-for-latency-iii-the-compiler:tut}

\begin{tutorial}
\textbf{Objectif.} Observer l'optimiseur agir sur un \omterm{def:ll:cpp-for-latency-iii-the-compiler:ub}{comportement indéfini}, le détecter avec un sanitiseur, et choisir les options par la mesure.
\textbf{État final~:} les trois listings d'assembleur, un rapport du sanitiseur, et
la \cref{fig:ll:cpp-for-latency-iii-the-compiler:flags}.
\begin{tutsteps}
\item \textbf{Trois contrats.} \texttt{python ll\_ub.py} compile \texttt{ll\_ub.cpp} en assembleur avec \texttt{-O0} et
\texttt{-O2} et écrit les listings ci-dessus.
\item \textbf{Détecter un dépassement.} \texttt{ll\_ub.ubsan(2147483647)} construit un appel de \texttt{increment(x)},
\texttt{return x + 1}, avec \texttt{-fsanitize=undefined}, et l'exécute~: «~runtime error~: signed integer overflow~: 2147483647 + 1 cannot
be represented in type int~».
\item \textbf{Construire avec un profil.} \texttt{firm.flagbench} construit un ensemble guidé par profil en trois étapes~: construction
instrumentée, exécution d'entraînement, reconstruction avec le profil.
\omcode{code/firm/flagbench/firm_flagbench.py}{38}{51}{Construction d'un ensemble d'options, guidé par profil ou non.}\label{lst:ll:cpp-for-latency-iii-the-compiler:build}
\item \textbf{Lancer la matrice} avec \texttt{python bench\_flags.py}~: sept ensembles sur le décodeur, trois sur la somme en virgule flottante, chacun
vérifié par rapport à la somme de contrôle de \texttt{-O2}.
\end{tutsteps}
\textbf{Pour aller plus loin.} Trouver lequel des passages qu'\texttt{-O3} ajoute à \texttt{-O2} rend le décodeur plus rapide
en les ajoutant un à un à \texttt{-O2}~; entraîner la construction guidée par profil sur seulement les cent premiers messages de l'échantillon
et mesurer à nouveau.
\end{tutorial}

\section{Construction~: la matrice d'options}\label{bld:ll:cpp-for-latency-iii-the-compiler:flagbench}

\begin{build}
\bfield{Rôle} Les options de publication de la société sont choisies, et rechoisies à chaque mise à niveau du compilateur, par la mesure avec un
contrôle de justesse, et jamais par habitude.
\bfield{Interface} Python \texttt{firm\_flagbench}~: \texttt{FlagSet(name, flags, pgo)},
\texttt{build(sources, fs, out\_dir, train\_args, includes)}, \texttt{run\_matrix(sources, flag\_sets, args, reference, repeats, cpus)} renvoyant
des lignes avec name, flags, checksum, time, ok et speed-up. Le benchmark affiche sa somme de contrôle sur sa première ligne et son temps
par unité sur la deuxième.
\bfield{Règles} Un ensemble dont la somme de contrôle diffère de celle de la référence est marqué not ok et n'est jamais recommandé~; les ensembles guidés par profil
sont entraînés avec les arguments propres au benchmark (ou une charge de travail d'entraînement distincte)~; chaque exécution est figée et répétée.
\bfield{Tests d'acceptation} \texttt{code/firm/flagbench/tests/}~: un petit programme sous \texttt{-O2}, \texttt{-O0} et une construction guidée par profil
(toutes ok), et sous \texttt{-ffast-math} (not ok, sa somme en virgule flottante diffère).
\bfield{Pour aller plus loin} Séparer les charges de travail d'entraînement et de mesure~; présenter des centiles au lieu du meilleur temps~;
rechercher automatiquement les passages de \texttt{-O3} responsables.
\end{build}

\begin{omsources}
\item J.~Corbet, «~Fun with NULL pointers, part 1~», LWN.net, 20 juillet 2009.
\item C.~Lattner, «~What every C programmer should know about undefined behavior~», blog LLVM, 2011.
\item manuel de GCC 11.4, «~Options that control optimization~» et «~Common function attributes~».
\end{omsources}

\section{Exercices}

\begin{exercise}[$\star$]\label{exo:ll:cpp-for-latency-iii-the-compiler:1}
D'après la \cref{fig:ll:cpp-for-latency-iii-the-compiler:flags}, combien coûte le décodeur par message à \texttt{-O0},
\texttt{-O2} et avec toutes les optimisations, et combien coûterait chacun pour 5 millions de messages par seconde~?
\end{exercise}

\begin{exercise}[$\star$]\label{exo:ll:cpp-for-latency-iii-the-compiler:2}
Lesquelles de ces expressions sont indéfinies~: \texttt{INT\_MAX + 1}~; \texttt{UINT\_MAX + 1u}~; la lecture d'un \texttt{uint32\_t} dans un tampon d'octets
avec \texttt{std::memcpy}~; sa lecture en convertissant l'adresse du tampon en \texttt{uint32\_t*}~?
\end{exercise}

\begin{exercise}[$\star$]\label{exo:ll:cpp-for-latency-iii-the-compiler:3}
Pourquoi \texttt{plus\_one\_greater} se compile-t-il en une constante même à \texttt{-O0}, et pourquoi le sanitiseur ne peut-il pas le détecter~?
\end{exercise}

\begin{exercise}[$\star\star$]\label{exo:ll:cpp-for-latency-iii-the-compiler:4}
La somme exacte de $1/k$ pour $k = 1$ à $2^{20}$ vaut environ 14.44. Laquelle des sommes en virgule flottante de la
\cref{fig:ll:cpp-for-latency-iii-the-compiler:flags} en est la plus proche, et cela rend-il \texttt{-ffast-math} sûr~?
\end{exercise}

\begin{exercise}[$\star\star$]\label{exo:ll:cpp-for-latency-iii-the-compiler:5}
Un décodeur lit un prix sur 32~bits avec \texttt{*reinterpret\_cast<const std::uint32\_t*>(buf + 16)}. Nommer deux règles
qu'il enfreint et écrire la ligne correcte.
\end{exercise}

\begin{exercise}[$\star\star$]\label{exo:ll:cpp-for-latency-iii-the-compiler:6}
Une construction guidée par profil a été entraînée sur le rejeu d'une journée calme, puis déployée le jour d'un rééquilibrage
d'indice. Que peut-il mal se passer~?
\end{exercise}

\begin{exercise}[$\star\star\star$]\label{exo:ll:cpp-for-latency-iii-the-compiler:7}
\emph{Programmation.} Ajouter à la matrice un jeu d'options \texttt{-O2 -fno-strict-aliasing} et compiler \texttt{alias} avec celui-ci. Que
change-t-il dans l'assembleur, et dans la vitesse du décodeur~?
\end{exercise}

\begin{exercise}[$\star\star\star$]\label{exo:ll:cpp-for-latency-iii-the-compiler:8}
\emph{Trouver la faille.} «~Nous compilons avec \texttt{-Ofast -march=native} sur le serveur de construction~; c'était 30~\% plus
rapide dans notre indice de référence.~»
\end{exercise}

\section{Problème~: la construction rapide qui était fausse}

\begin{problem}[{Problème du week-end --- revue des options de publication}]\label{pb:ll:cpp-for-latency-iii-the-compiler:1}
La publication d'une équipe utilise \texttt{-O2}. Un ingénieur propose
\texttt{-O3 -march=native -flto} avec \omterm{def:ll:cpp-for-latency-iii-the-compiler:pgo}{optimisation guidée par profil}, un autre \texttt{-Ofast}. L'équipe exécute \texttt{firm.flagbench}
(\cref{fig:ll:cpp-for-latency-iii-the-compiler:flags}).

\medskip\noindent\textbf{Partie I --- Le décodeur.}
\begin{enumerate}
\item Quel gain de vitesse par rapport à \texttt{-O2} chaque jeu d'options procure-t-il~?
\item Quel changement unique apporte la plus grande part du gain~?
\item Pourquoi l'\omterm{def:ll:cpp-for-latency-iii-the-compiler:lto}{optimisation à l'édition de liens} seule n'a-t-elle rien apporté ici~?
\item Un jeu d'options a-t-il modifié la somme de contrôle du décodeur~?
\end{enumerate}

\medskip\noindent\textbf{Partie II --- La somme en virgule flottante.}
\begin{enumerate}[resume]
\item Que donnent les trois jeux d'options, en temps et en valeur~?
\item Pourquoi \texttt{-ffast-math} a-t-il modifié la valeur~?
\item Quelle valeur est la plus proche de la somme exacte~?
\item Pourquoi l'équipe doit-elle malgré tout refuser \texttt{-ffast-math} pour son code de tarification et de risque~?
\end{enumerate}

\medskip\noindent\textbf{Partie III --- \omterm{def:ll:cpp-for-latency-iii-the-compiler:ub}{Comportement indéfini}.}
\begin{enumerate}[resume]
\item Le décodeur lit les champs par transtypage. Pourquoi cela peut-il fonctionner à \texttt{-O2} et casser à \texttt{-O3}~?
\item Que signale le sanitiseur de \omterm{def:ll:cpp-for-latency-iii-the-compiler:ub}{comportement indéfini} pour un incrément qui déborde~?
\item Dans quelles constructions les sanitiseurs doivent-ils être exécutés~?
\item Qu'enseigne le bogue du noyau de 2009 à une équipe qui vérifie les pointeurs après les avoir utilisés~?
\end{enumerate}

\medskip\noindent\textbf{Partie IV --- La décision.}
\begin{enumerate}[resume]
\item Énoncer le \emph{résultat nommé}~: le gain de vitesse des optimisations combinées par rapport à \texttt{-O2} sur le décodeur, et ce que
\texttt{-ffast-math} a fait à la somme en virgule flottante.
\item Quelles options la publication doit-elle utiliser, et qu'est-ce qui doit les accompagner~?
\item Comment choisir la journée d'entraînement du profil~?
\item À quoi \texttt{-march=native} engage-t-il la société~?
\item Comment le \omterm{def:ll:cpp-for-latency-iii-the-compiler:fmv}{multiversionnage de fonctions} modifierait-il cela~?
\item À quelle fréquence la matrice doit-elle être relancée~?
\item Pourquoi vérifier des propriétés de l'assembleur plutôt que les octets dans les tests~?
\item En une phrase~: que vaut une option de compilation sans somme de contrôle~?
\end{enumerate}
\end{problem}

\section{Questions d'entretien}

\begin{interviewq}[$\star$ \iqroles{développeur}]\label{iq:ll:cpp-for-latency-iii-the-compiler:1}
Qu'est-ce que le \omterm{def:ll:cpp-for-latency-iii-the-compiler:ub}{comportement indéfini}, et pourquoi existe-t-il en C++~?
\end{interviewq}

\begin{interviewq}[$\star\star$ \iqroles{développeur}]\label{iq:ll:cpp-for-latency-iii-the-compiler:2}
Qu'est-ce que la règle d'\omterm{def:ll:cpp-for-latency-iii-the-compiler:alias}{aliasing strict}~? Comment lisez-vous correctement un entier depuis un tampon réseau~?
\end{interviewq}

\begin{interviewq}[$\star\star$ \iqroles{développeur}]\label{iq:ll:cpp-for-latency-iii-the-compiler:3}
Que fait l'\omterm{def:ll:cpp-for-latency-iii-the-compiler:pgo}{optimisation guidée par profil}, et quels en sont les risques~?
\end{interviewq}

\begin{interviewq}[$\star\star$ \iqroles{développeur}]\label{iq:ll:cpp-for-latency-iii-the-compiler:4}
Utiliseriez-vous \texttt{-ffast-math} dans un système de trading~? Où, et où non~?
\end{interviewq}

\begin{interviewq}[$\star\star$ \iqroles{développeur}]\label{iq:ll:cpp-for-latency-iii-the-compiler:5}
Une vérification \texttt{if (ptr == nullptr)} disparaît du binaire optimisé. Pourquoi cela peut-il arriver~?
\end{interviewq}

\begin{interviewq}[$\star\star\star$ \iqroles{développeur}]\label{iq:ll:cpp-for-latency-iii-the-compiler:6}
Comment décideriez-vous, et continueriez-vous de décider, les options de compilation d'un binaire critique pour la latence~?
\end{interviewq}
